% GENERATED FILE. Edit canonical lessons, then run npm run generate:books.
% canonical-source-hash: fnv1a64:5a56d853abfd4a65
% canonical-lessons: RU-C276-desyatiletiye, RU-C276-period, RU-C276-etap, RU-C276-promezhutok, RU-C276-srok

\chapter{Spans of time}
\label{ch:ru-spans-of-time}
\hlchaptermodality{276}

\begin{chapteropening}
\textbf{By the end of this chapter:} \emph{I can talk about spans of time: a decade, a period, a stage, an interval and a deadline.}

Complete the last lesson of chapter 276: I can talk about spans of time: a decade, a period, a stage, an interval and a deadline.
\end{chapteropening}

\section[\texorpdfstring{\ru{десятилетие} (desyatilétiye)}{десятилетие (desyatilétiye)}]{\ru{десятилетие} (desyatilétiye) — a decade}

\label{lesson:RU-C276-desyatiletiye}



\begin{quote}
Before the new one: say the Russian for to happen again, then the Russian for to coincide.
\end{quote}

\subsection*{You'll want to know: \ru{десятилетие}}
\textbf{\ru{десятилетие}} (\emph{desyatilétiye}) — \textquotedblleft{}a decade\textquotedblright{}.

\begin{quote}
\ru{За} \ru{последнее} \ru{десятилетие} \ru{город} \ru{сильно} \ru{изменился}.
\end{quote}

\begin{quote}
\emph{Za poslédneye desyatilétiye górod sílno izmenílsya.}
\end{quote}

\begin{quote}
Over the last decade the town has changed a lot.
\end{quote}

Neuter. Ten years: \emph{\ru{десять} \ru{лет}}.

\begin{cousinweb}
\emph{\ru{десятилетие}} joins \emph{\ru{десять}} \textquotedblleft{}ten\textquotedblright{} and \emph{\ru{лето}} \textquotedblleft{}summer\textquotedblright{}, which is inherited from Proto-Slavic \emph{lěto} \textquotedblleft{}year\textquotedblright{} (EtymDB): ten years.
\end{cousinweb}

\begin{grammarlens}[title={What you've built}]
One more word for telling, explaining and describing at B1.
\end{grammarlens}

\subsection*{Your turn}
\begin{itemize}
\raggedright
  \item \emph{Say it:} \emph{desyatilétiye}
  \item \emph{Say it:} \emph{desyatilétiye}, once more
  \item \emph{Say it:} a sentence of your own with \emph{desyatilétiye}
  \item \emph{Recall:} say the Russian for to happen again, then the Russian for to coincide, then say \emph{desyatilétiye} again
\end{itemize}

\begin{tcolorbox}[breakable,colback=teal!4,colframe=teal!35!black,title={Before you move on}]
What does \ru{десятилетие} mean? (\textquotedblleft{}A decade\textquotedblright{}.) Say it once more.
\end{tcolorbox}
\section[\texorpdfstring{\ru{период} (períod)}{период (períod)}]{\ru{период} (períod) — a period, a time}

\label{lesson:RU-C276-period}



\begin{quote}
Before the new one: say the Russian for to coincide, then the Russian for a decade.
\end{quote}

\subsection*{You'll want to know: \ru{период}}
\textbf{\ru{период}} (\emph{períod}) — \textquotedblleft{}a period, a time\textquotedblright{}.

\begin{quote}
\ru{Это} \ru{был} \ru{трудный} \ru{период} \ru{в} \ru{его} \ru{жизни}.
\end{quote}

\begin{quote}
\emph{Éto byl trúdnyy períod v yevó zhízni.}
\end{quote}

\begin{quote}
It was a hard period in his life.
\end{quote}

Masculine. \emph{\ru{В} \ru{период}} + genitive — during: \emph{\ru{в} \ru{период} \ru{дождей}}.

\begin{cousinweb}
\emph{\ru{период}} is a borrowing usually traced through French or German to Greek \emph{períodos} \textquotedblleft{}a going round\textquotedblright{}, from \emph{perí} \textquotedblleft{}around\textquotedblright{} and \emph{hodós} \textquotedblleft{}way\textquotedblright{}; EtymDB lists the Russian word with no parent.
\end{cousinweb}

\begin{grammarlens}[title={What you've built}]
One more word for telling, explaining and describing at B1.
\end{grammarlens}

\subsection*{Your turn}
\begin{itemize}
\raggedright
  \item \emph{Say it:} \emph{períod}
  \item \emph{Say it:} \emph{períod}, once more
  \item \emph{Say it:} a sentence of your own with \emph{períod}
  \item \emph{Recall:} say the Russian for to coincide, then the Russian for a decade, then say \emph{períod} again
\end{itemize}

\begin{tcolorbox}[breakable,colback=teal!4,colframe=teal!35!black,title={Before you move on}]
What does \ru{период} mean? (\textquotedblleft{}A period, a time\textquotedblright{}.) Say it once more.
\end{tcolorbox}
\section[\texorpdfstring{\ru{этап} (etáp)}{этап (etáp)}]{\ru{этап} (etáp) — a stage, a phase}

\label{lesson:RU-C276-etap}



\begin{quote}
Before the new one: say the Russian for a decade, then the Russian for a period.
\end{quote}

\subsection*{You'll want to know: \ru{этап}}
\textbf{\ru{этап}} (\emph{etáp}) — \textquotedblleft{}a stage, a phase\textquotedblright{}.

\begin{quote}
\ru{Первый} \ru{этап} \ru{работы} \ru{закончен}.
\end{quote}

\begin{quote}
\emph{Pérvyy etáp rabóty zakónchen.}
\end{quote}

\begin{quote}
The first stage of the work is finished.
\end{quote}

Masculine. \emph{\ru{На} \ru{этом} \ru{этапе}} — at this stage.

\begin{cousinweb}
\emph{\ru{этап}} was borrowed from French \emph{étape} \textquotedblleft{}a stage, a stopping place\textquotedblright{} (EtymDB), from Old French \emph{estaple} \textquotedblleft{}a warehouse, a market\textquotedblright{}, itself from Middle Dutch \emph{stapel} \textquotedblleft{}a pillar, a market\textquotedblright{} (EtymDB).
\end{cousinweb}

\begin{grammarlens}[title={What you've built}]
One more word for telling, explaining and describing at B1.
\end{grammarlens}

\subsection*{Your turn}
\begin{itemize}
\raggedright
  \item \emph{Say it:} \emph{etáp}
  \item \emph{Say it:} \emph{etáp}, once more
  \item \emph{Say it:} a sentence of your own with \emph{etáp}
  \item \emph{Recall:} say the Russian for a decade, then the Russian for a period, then say \emph{etáp} again
\end{itemize}

\begin{tcolorbox}[breakable,colback=teal!4,colframe=teal!35!black,title={Before you move on}]
What does \ru{этап} mean? (\textquotedblleft{}A stage, a phase\textquotedblright{}.) Say it once more.
\end{tcolorbox}
\section[\texorpdfstring{\ru{промежуток} (promezhútok)}{промежуток (promezhútok)}]{\ru{промежуток} (promezhútok) — an interval, a gap}

\label{lesson:RU-C276-promezhutok}



\begin{quote}
Before the new one: say the Russian for a period, then the Russian for a stage.
\end{quote}

\subsection*{You'll want to know: \ru{промежуток}}
\textbf{\ru{промежуток}} (\emph{promezhútok}) — \textquotedblleft{}an interval, a gap\textquotedblright{}.

\begin{quote}
\ru{В} \ru{промежутке} \ru{между} \ru{уроками} \ru{мы} \ru{пили} \ru{чай}.
\end{quote}

\begin{quote}
\emph{V promezhútke mézhdu urókami my píli chay.}
\end{quote}

\begin{quote}
In the gap between lessons we drank tea.
\end{quote}

Genitive \emph{\ru{промежутка}}. \emph{\ru{Промежуток} \ru{времени}} — a stretch of time.

\begin{cousinweb}
\emph{\ru{промежуток}} is built on \emph{\ru{промеж}} \textquotedblleft{}between\textquotedblright{}, a cousin of \emph{\ru{между}}; both go back to the Slavic root of \emph{\ru{межа}} \textquotedblleft{}a boundary strip\textquotedblright{}, Proto-Slavic \emph{medja} (EtymDB). EtymDB lists no parent for \emph{\ru{промежуток}}.
\end{cousinweb}

\begin{grammarlens}[title={What you've built}]
One more word for telling, explaining and describing at B1.
\end{grammarlens}

\subsection*{Your turn}
\begin{itemize}
\raggedright
  \item \emph{Say it:} \emph{promezhútok}
  \item \emph{Say it:} \emph{promezhútok}, once more
  \item \emph{Say it:} a sentence of your own with \emph{promezhútok}
  \item \emph{Recall:} say the Russian for a period, then the Russian for a stage, then say \emph{promezhútok} again
\end{itemize}

\begin{tcolorbox}[breakable,colback=teal!4,colframe=teal!35!black,title={Before you move on}]
What does \ru{промежуток} mean? (\textquotedblleft{}An interval, a gap\textquotedblright{}.) Say it once more.
\end{tcolorbox}
\section[\texorpdfstring{\ru{срок} (srok)}{срок (srok)}]{\ru{срок} (srok) — a deadline, a time limit; a term}

\label{lesson:RU-C276-srok}



\begin{quote}
Before the new one: say the Russian for a stage, then the Russian for an interval.
\end{quote}

\subsection*{You'll want to know: \ru{срок}}
\textbf{\ru{срок}} (\emph{srok}) — \textquotedblleft{}a deadline, a time limit; a term\textquotedblright{}.

\begin{quote}
\ru{Срок} \ru{сдачи} \ru{проекта} — \ru{пятница}.
\end{quote}

\begin{quote}
\emph{Srok sdáchi proyékta — pyátnitsa.}
\end{quote}

\begin{quote}
The deadline for the project is Friday.
\end{quote}

\emph{\ru{В} \ru{срок}} — on time; \emph{\ru{срок} \ru{годности}} — shelf life.

\begin{cousinweb}
\emph{\ru{срок}} is usually traced to \emph{\ru{с}-} and the root of \emph{\ru{речь}} \textquotedblleft{}speech\textquotedblright{}, Proto-Slavic \emph{rekti} \textquotedblleft{}to say\textquotedblright{} (EtymDB lists \emph{\ru{речь}}), so \textquotedblleft{}what was agreed\textquotedblright{}; EtymDB lists no parent, and the origin is not fully settled.
\end{cousinweb}

\begin{grammarlens}[title={What you've built}]
One more word for telling, explaining and describing at B1.
\end{grammarlens}

\subsection*{Your turn}
\begin{itemize}
\raggedright
  \item \emph{Say it:} \emph{srok}
  \item \emph{Say it:} \emph{srok}, once more
  \item \emph{Say it:} a sentence of your own with \emph{srok}
  \item \emph{Recall:} say the Russian for a stage, then the Russian for an interval, then say \emph{srok} again
\end{itemize}

\begin{tcolorbox}[breakable,colback=teal!4,colframe=teal!35!black,title={Before you move on}]
What does \ru{срок} mean? (\textquotedblleft{}A deadline, a time limit; a term\textquotedblright{}.) Say it once more.
\end{tcolorbox}
